%% kapa: exam=1 round=24 year=2013 period=1 form=A subject=civil_law subject_ko=민법 seq=29 number=29
%% source: Q-Net 감정평가사 기출문제 (cst00309) artlSeq=5073802 file=extracted/1008/5073802/2013년 제24회 감정평가사 1차시험 기출문제/1교시(A형).hwp
%% conversion: kordoc 4.24.51 -> tools/gen1.py (auto); review=proofread-sonnet; flags=none
\begin{question}{29}
점유에 관한 설명으로 옳지 않은 것은? (다툼이 있으면 판례에 의함)
\begin{choices}
  \item 피상속인의 사망을 알지 못한 상속인은 피상속인의 점유를 승계하지 못한다.
  \item 점유자는 과실 없이 점유하는 것으로 추정되지 않는다.
  \item 점유물반환청구에 대하여 점유침탈자는 점유물에 대한 본권이 있다는 이유로 반환을 거부할 수 없다.
  \item 점유자가 상대방의 사기로 물건을 인도한 경우에는 점유물반환청구를 할 수 없다.
  \item 직접점유가 침탈된 경우 직접점유자가 반환받을 수 없는 때에는 간접점유자는 자기에게 반환할 것을 청구할 수 있다.
\end{choices}
\end{question}
\answer{1}
